\documentclass[a4paper,12pt]{article}

% ==================== ПАКЕТЫ ====================
\usepackage[T2A]{fontenc}
\usepackage[utf8]{inputenc}
\usepackage[russian]{babel}
\usepackage{amsmath,amssymb}
\usepackage{geometry}
\geometry{top=2cm,bottom=2cm,left=2.5cm,right=1.5cm}
\usepackage{graphicx, multirow}
\usepackage{booktabs}
\usepackage{array}
\usepackage{float}
\usepackage{enumitem}
\usepackage{indentfirst}
\usepackage{caption}
\usepackage{pgfplots}
\pgfplotsset{compat=1.18}
\usepackage[hidelinks]{hyperref}

% ==================== НАСТРОЙКИ ОФОРМЛЕНИЯ ====================
\captionsetup{
    font=small,
    labelfont=bf,
    justification=centering,
    skip=6pt
}
\setlist{nosep, leftmargin=1.6em}
\linespread{1.06}
\setlength{\parindent}{1.25em}

\pgfplotsset{
    every axis/.append style={
        grid=major,
        grid style={line width=.1pt, draw=gray!25},
        tick label style={font=\small},
        label style={font=\small},
        legend style={
            font=\footnotesize,
            fill=white, fill opacity=0.9, text opacity=1,
            draw=gray!40,
            cells={anchor=west}
        }
    }
}

\title{Эффект Холла в полупроводниках}
\author{}
\date{\today}

% ==================== ДОКУМЕНТ ====================
\begin{document}

% ==================== ТИТУЛ ====================
\begin{center}
    {\Large\bfseries Эффект Холла в полупроводниках\par}
    \vspace{0.4cm}
    {\large Имя Фамилия, группа ХХХ\par}
    \vspace{0.2cm}
    {\large \today\par}
\end{center}

\vspace{0.5cm}

% ==================== 1. АННОТАЦИЯ ====================
\section{Аннотация}
Целью данной работы является измерение подвижности и концентрации носителей заряда в полупроводниках. Исследуется эффект Холла в образце легированного германия: измеряется ЭДС Холла при различных значениях магнитной индукции и тока через образец, по угловому коэффициенту зависимости определяется постоянная Холла, знак которой позволяет установить тип проводимости. Дополнительно измеряется удельная проводимость образца, что даёт возможность вычислить подвижность носителей.

% ==================== 2. ТЕОРЕТИЧЕСКОЕ ВВЕДЕНИЕ ====================
\section{Теоретическое введение}

\subsection{Сила Лоренца и ЭДС Холла}
На электрон, движущийся в магнитном поле, действует сила Лоренца. Если пластину с током поместить в магнитное поле, на носителях заряда возникает дополнительная сила, приводящая к появлению поперечной разности потенциалов. Полная сила, действующая на электрон:
\begin{equation}
    F_1 = -eE - e\langle v \rangle B .
\end{equation}

Под действием этой силы электроны отклоняются к грани Б, а на грани А создаётся нескомпенсированный положительный заряд. Из-за разности потенциалов возникает электрическое поле, направленное от грани А к грани Б: $F_2 = eE_z$. Приравнивая $F_1$ и $F_2$, найдём ЭДС Холла:
\begin{equation}
    U_{ab} = -\frac{IB}{nea} = -R_x \frac{IB}{a},
    \qquad R_x = \frac{1}{ne}.
\end{equation}
Величина $R_x$ называется постоянной Холла. Её знак определяется знаком носителей заряда: для электронной проводимости $R_x < 0$, для дырочной — $R_x > 0$.

\subsection{Удельная проводимость и подвижность}
Удельная проводимость образца измеряется по падению напряжения на участке известной длины:
\begin{equation}
    \sigma = \frac{I\,L_{34}}{U_{35}\,a\,l}.
\end{equation}

Зная постоянную Холла и удельную проводимость, можно найти подвижность носителей:
\begin{equation}
    b = \sigma\,|R_x| .
\end{equation}

% ==================== 3. ОПИСАНИЕ УСТАНОВКИ ====================
\section{Описание установки}
Исследования проводятся на лабораторной установке для изучения эффекта Холла. Схема установки включает источник питания магнита, цифровой вольтметр В7-78/1, блок управления и образец германия с разъёмами.

В работе используются:
\begin{itemize}
    \item электромагнит с источником питания;
    \item амперметр;
    \item миллиамперметр;
    \item милливеберметр;
    \item реостат;
    \item источник питания ($1{,}5$ В);
    \item образцы легированного германия.
\end{itemize}

Параметры исследуемого образца: $a = 2{,}2$ мм (толщина), $l = 7$ мм и $L_{34} = 3$ мм (характерные размеры между контактами).

% ==================== 4. РЕЗУЛЬТАТЫ И ОБСУЖДЕНИЯ ====================
\section{Результаты и обсуждения}

\subsection{Калибровка электромагнита}
Для определения связи между индукцией магнитного поля $B$ в зазоре электромагнита и током $I_M$ через обмотку проведена калибровка. Магнитный поток измеряется милливеберметром: $\Phi = B S N$, откуда индукция $B = \Phi / S_{\mathrm{eff}}$. Результаты приведены в таблице~\ref{tab:calib} и на рисунке~\ref{fig:calib}. Погрешность милливеберметра рассчитана по формуле

\begin{equation}
	\delta\Phi = 0,005\Phi + 2 \text{мВб}
\end{equation}

Погрешность для амперметра на источника положена $\delta I_M = 0,03$ А
\begin{table}[H]
    \centering
    \caption{Калибровка электромагнита}
    \label{tab:calib}
    \resizebox{\textwidth}{!}{%
        \begin{tabular}{l *{17}{c}}
            \toprule
            $I_M$, А    & 0.13 & 0.26 & 0.39 & 0.52 & 0.65 & 0.78 & 0.91 & 1.04 & 1.17 & 1.30 & 1.43 & 1.56 & 1.69 & 1.82 & 1.95 & 2.08 & 2.10 \\
            $\Phi$, мВб & 7    & 11   & 18   & 24   & 29   & 35   & 40   & 45   & 50   & 54   & 57   & 60   & 63   & 65   & 67   & 69   & 69   \\
            $B$, Тл     & 0.100 & 0.157 & 0.257 & 0.343 & 0.414 & 0.500 & 0.571 & 0.643 & 0.714 & 0.771 & 0.814 & 0.857 & 0.900 & 0.929 & 0.957 & 0.986 & 0.986 \\
            \bottomrule
        \end{tabular}%
    }
\end{table}

\begin{figure}[H]
    \centering
    \begin{tikzpicture}
        \begin{axis}[
            width=0.85\textwidth, height=6.5cm,
            xlabel={$I_M$, А}, ylabel={$B$, Тл},
            legend pos=north west,
            xmin=0, xmax=2.2, ymin=0
        ]
            \addplot[color=blue, mark=*, thick] coordinates {
                (0.13, 0.100) (0.26, 0.157) (0.39, 0.257) (0.52, 0.343)
                (0.65, 0.414) (0.78, 0.500) (0.91, 0.571) (1.04, 0.643)
                (1.17, 0.714) (1.30, 0.771) (1.43, 0.814) (1.56, 0.857)
                (1.69, 0.900) (1.82, 0.929) (1.95, 0.957) (2.08, 0.986) (2.10, 0.986)
            };
            \legend{Калибровка $B(I_M)$}
        \end{axis}
    \end{tikzpicture}
    \caption{График калибровки электромагнита}
    \label{fig:calib}
\end{figure}

\subsection{Измерение ЭДС Холла}
Снимем зависимость напряжения $|U_{34}|$ от магнитной индукции $B$. При $I_M = 0$ напряжение не равно нулю: $U_{034} = 0{,}360$ мВ, поэтому далее рассматривается $|U_{34}|$ с учётом этой поправки. Серия экспериментов проведена для различных токов через образец $I$ (от $0{,}2$ до $1{,}1$ мА). Все измеренные значения напряжений имели отрицательный знак. Результаты приведены в таблице~\ref{tab:hall}, соответствующие зависимости — на рисунке~\ref{fig:hall}.

\begin{table}[H]
    \centering
    \caption{Зависимость напряжения $|U_{34}|$ (мВ) от индукции поля $B$ (Тл)}
    \label{tab:hall}
    \small
    \setlength{\tabcolsep}{5pt}
    \begin{tabular}{c|*{7}{c}}
        \toprule
        \multirow{2}{*}{$B$, Тл} & \multicolumn{7}{c}{$|U_{34}|$, мВ, при токе $I$, мА} \\
        \cmidrule(lr){2-8}
         & 0.20 & 0.40 & 0.55 & 0.70 & 0.90 & 1.10 & 1.10\textsuperscript{*} \\
        \midrule
        0.000 & 0.069 & 0.115 & 0.157 & 0.200 & 0.261 & 0.325 & 0.325 \\
        0.100 & 0.339 & 0.719 & 0.999 & 1.263 & 1.568 & 1.822 & 1.815 \\
        0.257 & 0.944 & 1.828 & 2.538 & 3.222 & 4.136 & 5.152 & 5.062 \\
        0.414 & 1.513 & 2.971 & 4.115 & 5.231 & 6.751 & 8.143 & 8.142 \\
        0.571 & 2.027 & 3.997 & 5.569 & 6.977 & 9.056 & 10.972 & 10.931 \\
        0.714 & 2.473 & 4.871 & 6.722 & 8.515 & 10.991 & 13.342 & 13.365 \\
        0.814 & 2.825 & 5.532 & 7.598 & 9.682 & 12.495 & 15.182 & 15.180 \\
        0.900 & 3.053 & 5.987 & 8.208 & 10.441 & 13.485 & 16.419 & 16.215 \\
        0.957 & 3.225 & 6.298 & 8.687 & 11.029 & 14.244 & 17.317 & 17.165 \\
        \bottomrule
    \end{tabular}
    \par\vspace{2pt}
    {\footnotesize\textsuperscript{*} повторное измерение при $I = 1{,}10$ мА}
\end{table}

\begin{figure}[H]
    \centering
    \begin{tikzpicture}
        \begin{axis}[
            width=0.92\textwidth, height=9.5cm,
            xlabel={$B$, Тл}, ylabel={$|U_{34}|$, мВ},
            xmin=0, xmax=1.0, ymin=0,
            legend columns=2,
            legend style={at={(0.02,0.98)}, anchor=north west}
        ]
            \addplot[only marks, mark=*, color=blue] coordinates {
                (0.000,0.069) (0.100,0.339) (0.257,0.944) (0.414,1.513) (0.571,2.027) (0.714,2.473) (0.814,2.825) (0.900,3.053) (0.957,3.225)};
            \addplot[domain=0:1, color=blue, thick, no marks] {3.30*x + 0.10};
            \addlegendentry{$I=0.20$ мА}

            \addplot[only marks, mark=square*, color=red] coordinates {
                (0.000,0.115) (0.100,0.719) (0.257,1.828) (0.414,2.971) (0.571,3.997) (0.714,4.871) (0.814,5.532) (0.900,5.987) (0.957,6.298)};
            \addplot[domain=0:1, color=red, thick, no marks] {6.48*x + 0.15};
            \addlegendentry{$I=0.40$ мА}

            \addplot[only marks, mark=triangle*, color=brown] coordinates {
                (0.000,0.157) (0.100,0.999) (0.257,2.538) (0.414,4.115) (0.571,5.569) (0.714,6.722) (0.814,7.598) (0.900,8.208) (0.957,8.687)};
            \addplot[domain=0:1, color=brown, thick, no marks] {8.94*x + 0.18};
            \addlegendentry{$I=0.55$ мА}

            \addplot[only marks, mark=diamond*, color=green!60!black] coordinates {
                (0.000,0.200) (0.100,1.263) (0.257,3.222) (0.414,5.231) (0.571,6.977) (0.714,8.515) (0.814,9.682) (0.900,10.441) (0.957,11.029)};
            \addplot[domain=0:1, color=green!60!black, thick, no marks] {11.35*x + 0.22};
            \addlegendentry{$I=0.70$ мА}

            \addplot[only marks, mark=pentagon*, color=orange] coordinates {
                (0.000,0.261) (0.100,1.568) (0.257,4.136) (0.414,6.751) (0.571,9.056) (0.714,10.991) (0.814,12.495) (0.900,13.485) (0.957,14.244)};
            \addplot[domain=0:1, color=orange, thick, no marks] {14.65*x + 0.28};
            \addlegendentry{$I=0.90$ мА}

            \addplot[only marks, mark=star, color=magenta] coordinates {
                (0.000,0.325) (0.100,1.822) (0.257,5.152) (0.414,8.143) (0.571,10.972) (0.714,13.342) (0.814,15.182) (0.900,16.419) (0.957,17.317)};
            \addplot[domain=0:1, color=magenta, thick, no marks] {17.85*x + 0.35};
            \addlegendentry{$I=1.10$ мА}

            \addplot[only marks, mark=x, color=gray] coordinates {
                (0.000,0.325) (0.100,1.815) (0.257,5.062) (0.414,8.142) (0.571,10.931) (0.714,13.365) (0.814,15.180) (0.900,16.215) (0.957,17.165)};
            \addplot[domain=0:1, color=gray, dashed, thick, no marks] {17.65*x + 0.35};
            \addlegendentry{$I=1.10$ мА (повт.)}
        \end{axis}
    \end{tikzpicture}
    \caption{Семейство зависимостей ЭДС Холла от индукции магнитного поля при различных токах через образец}
    \label{fig:hall}
\end{figure}

Из рисунка~\ref{fig:hall} видно, что при фиксированном токе $I$ напряжение $|U_{34}|$ линейно растёт с индукцией $B$, а наклон прямых увеличивается пропорционально току. Это качественно согласуется с формулой (2).

\subsection{Определение постоянной Холла}
Угловые коэффициенты аппроксимирующих прямых $k = \mathrm{d}U/\mathrm{d}B$ (рис.~\ref{fig:hall}) сведены в таблицу~\ref{tab:k}. График зависимости $k = f(I)$ приведён на рисунке~\ref{fig:k}.

\begin{table}[H]
    \centering
    \caption{Угловые коэффициенты аппроксимирующих прямых}
    \label{tab:k}
    \begin{tabular}{c ccccc}
        \toprule
        $I$, мА   & 0.20 & 0.40 & 0.55 & 0.70 & 0.90 \\
        \midrule
        $k$, мВ/Тл & 3.30 & 6.48 & 8.94 & 11.35 & 14.65 \\
        \bottomrule
    \end{tabular}
    \qquad
    \begin{tabular}{c c}
        \toprule
        $I$, мА & $k$, мВ/Тл \\
        \midrule
        1.10 & 17.85 \\
        1.10\textsuperscript{*} & 17.65 \\
        \bottomrule
    \end{tabular}
    \par\vspace{2pt}
    {\footnotesize\textsuperscript{*} повторное измерение}
\end{table}

\begin{figure}[H]
    \centering
    \begin{tikzpicture}
        \begin{axis}[
            width=0.75\textwidth, height=6cm,
            xlabel={$I$, мА}, ylabel={$k$, мВ/Тл},
            legend pos=north west,
            xmin=0, xmax=1.2, ymin=0
        ]
            \addplot[only marks, mark=*, color=blue] coordinates {
                (0.20, 3.30) (0.40, 6.48) (0.55, 8.94) (0.70, 11.35) (0.90, 14.65) (1.10, 17.85)
            };
            \addplot[domain=0:1.2, color=red, thick, no marks] {16.15*x + 0.05};
            \legend{Эксперимент, $k(I) = 16{,}15 \cdot I$}
        \end{axis}
    \end{tikzpicture}
    \caption{Зависимость углового коэффициента от тока}
    \label{fig:k}
\end{figure}

По графику найдём постоянную Холла $R_x$. Угловой коэффициент прямой $K = \mathrm{d}k/\mathrm{d}I = 16{,}15$ В/(Тл$\cdot$А). Погрешность рассчитана по методу наименьших квадратов:
\begin{equation}
    R_x = -K a = -16{,}15 \cdot 2{,}2\cdot 10^{-3}
        = -(3{,}55 \pm 0{,}24)\cdot 10^{-2}\ \text{м}^3/\text{Кл}.
\end{equation}
Относительная погрешность составляет $6{,}8\%$.

\subsection{Определение типа проводимости}
С учётом направления тока в образце и отрицательного знака ЭДС Холла по правилу векторного произведения получаем, что проводимость образца — \textbf{электронная} ($n$-тип).

\subsection{Концентрация, проводимость и подвижность носителей}
Концентрация носителей тока рассчитывается как
\begin{equation}
    n = \frac{1}{|R_x|\,e} = (1{,}76 \pm 0{,}12)\cdot 10^{20}\ \text{м}^{-3}.
\end{equation}

По формуле (3) рассчитаем удельную проводимость образца. При $I = 1{,}1$ мА, $U_{35} = 0{,}345$ В и параметрах установки $a = 2{,}2$ мм, $l = 7$ мм, $L_{34} = 3$ мм:
\begin{equation}
    \sigma = \frac{I\,L_{34}}{U_{35}\,a\,l}
           = 8{,}70\ \text{(Ом}\cdot\text{м)}^{-1}.
\end{equation}

Подвижность носителей в образце:
\begin{equation}
    b = \sigma\,|R_x| = (3088 \pm 260)\ \text{см}^2/(\text{В}\cdot\text{с}).
\end{equation}
Теоретическое значение для электронной проводимости чистого германия: $b_{\text{теор}} = 3800$ см$^2$/(В$\cdot$с).

% ==================== 5. ВЫВОДЫ ====================
\section{Выводы}
В ходе работы был исследован эффект Холла в полупроводнике германия $n$-типа. Определены следующие характеристики образца:
\begin{itemize}
    \item постоянная Холла $R_x = -3{,}55\cdot 10^{-2}$ м$^3$/Кл;
    \item концентрация электронов $n = 1{,}76\cdot 10^{20}$ м$^{-3}$;
    \item удельная электрическая проводимость $\sigma = 8{,}70$ (Ом$\cdot$м)$^{-1}$;
    \item подвижность носителей $b \approx 3090$ см$^2$/(В$\cdot$с).
\end{itemize}

Результаты совпадают с табличными по порядку величины. Небольшое отличие экспериментального значения подвижности от теоретического объясняется наличием примесей в легированном германии, которые снижают подвижность за счёт дополнительного рассеяния носителей заряда.

\end{document}